% =====================================================================
%  LinkStudio manual - preamble (packages, styles, macros)
% =====================================================================
\PassOptionsToPackage{obeyspaces,spaces}{url}
\usepackage[T1]{fontenc}
\usepackage[utf8]{inputenc}
\usepackage{lmodern}
\usepackage[a4paper,top=27mm,bottom=27mm,inner=28mm,outer=24mm,headheight=14pt]{geometry}
\usepackage{microtype}
\usepackage{amsmath,amssymb,amsthm,bm,mathtools}
\usepackage{siunitx}
\sisetup{per-mode=symbol,group-separator={\,},range-phrase=--}
\usepackage{booktabs,longtable,array,tabularx,multirow}
\usepackage{enumitem}
\setlist{itemsep=2pt,topsep=4pt}
\usepackage{xcolor}
\usepackage{graphicx}
\usepackage{caption,subcaption}
\captionsetup{font=small,labelfont=bf,format=hang,justification=raggedright,singlelinecheck=true}
\usepackage{fancyhdr}
\usepackage{titlesec}
\usepackage{tikz}
\usetikzlibrary{arrows.meta,positioning,calc,shapes.geometric,shapes.misc,fit,backgrounds,
  decorations.pathreplacing,decorations.markings,matrix,patterns}
\usepackage{pgfplots}
\pgfplotsset{compat=1.18}
\usepgfplotslibrary{groupplots,fillbetween}
\usepackage{listings}
\usepackage[most]{tcolorbox}
\usepackage[hidelinks]{hyperref}
\usepackage[nameinlink,capitalise,noabbrev]{cleveref}

% ------------------------------------------------------------- colours
\definecolor{ink}{HTML}{1F2933}
\definecolor{accent}{HTML}{1D4E89}
\definecolor{accent2}{HTML}{C8553D}
\definecolor{accent3}{HTML}{2A9D8F}
\definecolor{accent4}{HTML}{8E5DB5}
\definecolor{accent5}{HTML}{E9A23B}
\definecolor{soft}{HTML}{EEF3F9}
\definecolor{softwarm}{HTML}{FBF1EC}
\definecolor{softgreen}{HTML}{EAF6F4}
\definecolor{codebg}{HTML}{F7F8FA}
\definecolor{codeframe}{HTML}{D5DBE3}
\definecolor{kwcolor}{HTML}{1D4E89}
\definecolor{strcolor}{HTML}{A23E48}
\definecolor{commentcolor}{HTML}{6B7785}
% signal-kind colours (same as the GUI)
\definecolor{kbits}{HTML}{8E5DB5}
\definecolor{kllr}{HTML}{D6589F}
\definecolor{ksym}{HTML}{2A9D8F}
\definecolor{ksamp}{HTML}{1D6FB8}
\definecolor{kgrid}{HTML}{E9A23B}
\definecolor{kcsi}{HTML}{C8553D}

% ------------------------------------------------------------- headings
\titleformat{\chapter}[display]
  {\normalfont\sffamily\color{accent}}
  {\large\scshape Chapter \thechapter}{4pt}
  {\Huge\bfseries}[\vspace{2pt}{\color{accent}\titlerule[0.8pt]}]
\titlespacing*{\chapter}{0pt}{-10pt}{24pt}
\titleformat{\part}[display]
  {\normalfont\sffamily\centering\color{accent}}
  {\Large\scshape Part \thepart}{8pt}
  {\Huge\bfseries}
\titleformat{\section}{\normalfont\sffamily\Large\bfseries\color{ink}}{\color{accent}\thesection}{0.8em}{}
\titleformat{\subsection}{\normalfont\sffamily\large\bfseries\color{ink}}{\color{accent}\thesubsection}{0.8em}{}
\titleformat{\subsubsection}{\normalfont\sffamily\normalsize\bfseries\color{ink}}{}{0em}{}

\pagestyle{fancy}
\fancyhf{}
\fancyhead[LE]{\sffamily\small\color{commentcolor}\leftmark}
\fancyhead[RO]{\sffamily\small\color{commentcolor}\rightmark}
\fancyfoot[C]{\sffamily\small\thepage}
\renewcommand{\headrulewidth}{0.3pt}

% ------------------------------------------------------------- theorems
\theoremstyle{definition}
\newtheorem{definition}{Definition}[chapter]
\newtheorem{proposition}{Proposition}[chapter]
\newtheorem{example}{Example}[chapter]
\theoremstyle{remark}
\newtheorem*{remark}{Remark}

% ------------------------------------------------------------- code
\lstdefinestyle{py}{
  language=Python,
  basicstyle=\ttfamily\footnotesize\color{ink},
  keywordstyle=\color{kwcolor}\bfseries,
  stringstyle=\color{strcolor},
  commentstyle=\color{commentcolor}\itshape,
  numbers=left, numberstyle=\tiny\color{commentcolor}, numbersep=6pt,
  backgroundcolor=\color{codebg},
  frame=single, rulecolor=\color{codeframe}, framesep=3pt,
  xleftmargin=16pt, framexleftmargin=14pt,
  breaklines=true, breakatwhitespace=false, columns=fullflexible, keepspaces=true,
  showstringspaces=false, tabsize=4, upquote=true,
  captionpos=t, abovecaptionskip=2pt, belowcaptionskip=4pt,
  morekeywords={self,None,True,False},
  literate={π}{{$\pi$}}1 {β}{{$\beta$}}1 {—}{{--}}1
}
\lstdefinestyle{yaml}{
  basicstyle=\ttfamily\footnotesize\color{ink},
  backgroundcolor=\color{codebg}, frame=single, rulecolor=\color{codeframe},
  breaklines=true, columns=fullflexible, keepspaces=true, showstringspaces=false,
  commentstyle=\color{commentcolor}\itshape, morecomment=[l]{\#},
  captionpos=t
}
\lstdefinestyle{shell}{
  basicstyle=\ttfamily\footnotesize\color{ink},
  backgroundcolor=\color{codebg}, frame=single, rulecolor=\color{codeframe},
  columns=fullflexible, keepspaces=true, showstringspaces=false, captionpos=t
}
\lstset{style=py}
\renewcommand{\lstlistingname}{Code}
\renewcommand{\lstlistlistingname}{List of code excerpts}
\crefname{lstlisting}{Code}{Code}
\Crefname{lstlisting}{Code}{Code}

% Excerpts are extracted from the code base by tools/generate.py.
% \snipdef{name}{file}{first}{last} is written to generated/snippets.tex.
\makeatletter
\newcommand{\snipdef}[4]{%
  \expandafter\xdef\csname snipfile@#1\endcsname{\detokenize{#2}}%
  \expandafter\gdef\csname snipfirst@#1\endcsname{#3}%
  \expandafter\gdef\csname sniplast@#1\endcsname{#4}}
\newcommand{\snipsource}[1]{%
  \texttt{\csname snipfile@#1\endcsname}, lines
  \csname snipfirst@#1\endcsname--\csname sniplast@#1\endcsname}
% \codeexcerpt{name}{caption}
\newcommand{\codeexcerpt}[2]{%
  \edef\snip@first{\csname snipfirst@#1\endcsname}%
  \lstinputlisting[style=py,firstnumber=\snip@first,label={lst:#1},
    caption={#2\hfill{\normalfont\footnotesize\color{commentcolor}\snipsource{#1}}}]
    {generated/snippets/#1.py}}
\makeatother
% Generated by manual/tools/generate.py: \snipdef{name}{file}{first line}{last line}
\snipdef{block_reset}{linkstudio/core/block.py}{102}{114}
\snipdef{graph_toposort}{linkstudio/core/graph.py}{193}{215}
\snipdef{graph_execute}{linkstudio/core/graph.py}{240}{265}
\snipdef{stop_check}{linkstudio/core/simulation.py}{35}{47}
\snipdef{parallel_merge}{linkstudio/core/parallel.py}{140}{145}
\snipdef{sweep_point}{linkstudio/core/sweep.py}{140}{146}
\snipdef{awgn_noise_var}{linkstudio/phy/channel/awgn.py}{49}{72}
\snipdef{awgn_process}{linkstudio/phy/channel/awgn.py}{74}{82}
\snipdef{wilson}{linkstudio/phy/metrics.py}{13}{21}
\snipdef{demap_soft}{linkstudio/phy/dsp/constellations.py}{72}{86}
\snipdef{qam_build}{linkstudio/phy/dsp/constellations.py}{118}{125}
\snipdef{rrc_taps}{linkstudio/phy/dsp/filters.py}{19}{36}
\snipdef{mf_normalize}{linkstudio/phy/waveform/pulse_shaping.py}{96}{106}
\snipdef{pulse_shaper}{linkstudio/phy/waveform/pulse_shaping.py}{46}{53}
\snipdef{phase_noise}{linkstudio/phy/channel/impairments.py}{74}{79}
\snipdef{iq_imbalance}{linkstudio/phy/channel/impairments.py}{117}{124}
\snipdef{power_amplifier}{linkstudio/phy/channel/impairments.py}{150}{163}
\snipdef{quantizer}{linkstudio/phy/channel/impairments.py}{191}{198}
\snipdef{sos}{linkstudio/phy/channel/fading.py}{68}{86}
\snipdef{flat_fading}{linkstudio/phy/channel/fading.py}{115}{126}
\snipdef{linear_eq}{linkstudio/phy/receiver.py}{87}{108}
\snipdef{lms_loop}{linkstudio/phy/receiver.py}{147}{166}
\snipdef{loop_gains}{linkstudio/phy/sync.py}{24}{31}
\snipdef{cubic}{linkstudio/phy/sync.py}{36}{44}
\snipdef{om_estimate}{linkstudio/phy/sync.py}{117}{123}
\snipdef{gardner_loop}{linkstudio/phy/sync.py}{128}{136}
\snipdef{frame_sync}{linkstudio/phy/sync.py}{191}{212}
\snipdef{dd_pll}{linkstudio/phy/sync.py}{246}{252}
\snipdef{ofdm_mod}{linkstudio/phy/waveform/ofdm.py}{94}{108}
\snipdef{ofdm_sync}{linkstudio/phy/waveform/ofdm.py}{148}{159}
\snipdef{ls_slope}{linkstudio/phy/waveform/ofdm.py}{235}{251}
\snipdef{mimo_linear}{linkstudio/phy/mimo.py}{152}{161}
\snipdef{alamouti_dec}{linkstudio/phy/mimo.py}{196}{200}
\snipdef{crc_matrix}{linkstudio/phy/coding/crc.py}{30}{39}
\snipdef{lfsr}{linkstudio/phy/dsp/lfsr.py}{27}{40}
\snipdef{viterbi}{linkstudio/phy/coding/convolutional.py}{67}{86}
\snipdef{ldpc_encode}{linkstudio/phy/coding/ldpc.py}{71}{74}
\snipdef{ldpc_minsum}{linkstudio/phy/coding/ldpc.py}{92}{103}
\snipdef{polar_rel}{linkstudio/phy/coding/polar.py}{35}{39}
\snipdef{polar_encode}{linkstudio/phy/coding/polar.py}{44}{52}
\snipdef{polar_sc}{linkstudio/phy/coding/polar.py}{59}{83}
\snipdef{depuncture}{linkstudio/phy/coding/puncture.py}{55}{64}
\snipdef{cable}{linkstudio/phy/wireline.py}{296}{305}
\snipdef{xtalk}{linkstudio/phy/wireline.py}{357}{362}
\snipdef{echo_nlms}{linkstudio/phy/wireline.py}{479}{484}
\snipdef{dfe_loop}{linkstudio/phy/wireline.py}{565}{578}
\snipdef{waterfill}{linkstudio/phy/wireline.py}{616}{619}
\snipdef{dmt_hermitian}{linkstudio/phy/wireline.py}{686}{692}
\snipdef{jitter}{linkstudio/phy/wireline.py}{602}{606}
\snipdef{wiener}{linkstudio/phy/optical.py}{51}{53}
\snipdef{im_mod}{linkstudio/phy/optical.py}{172}{176}
\snipdef{iq_mzm}{linkstudio/phy/optical.py}{280}{289}
\snipdef{disp_op}{linkstudio/phy/optical.py}{313}{324}
\snipdef{beta2}{linkstudio/phy/optical.py}{298}{302}
\snipdef{ssfm}{linkstudio/phy/optical.py}{333}{373}
\snipdef{pmd_transfer}{linkstudio/phy/optical.py}{426}{436}
\snipdef{ase}{linkstudio/phy/optical.py}{521}{530}
\snipdef{add_ase}{linkstudio/phy/optical.py}{533}{535}
\snipdef{osnr}{linkstudio/phy/optical.py}{68}{74}
\snipdef{link_loop}{linkstudio/phy/optical.py}{687}{696}
\snipdef{photodiode}{linkstudio/phy/optical.py}{1005}{1010}
\snipdef{coherent_rx}{linkstudio/phy/optical.py}{1054}{1073}
\snipdef{cd_comp}{linkstudio/phy/optical.py}{1140}{1141}
\snipdef{fo_est}{linkstudio/phy/optical.py}{1180}{1187}
\snipdef{butterfly}{linkstudio/phy/optical.py}{1251}{1265}
\snipdef{bps}{linkstudio/phy/optical.py}{1306}{1319}
\snipdef{wdm_demux}{linkstudio/phy/optical.py}{794}{800}
\snipdef{fso}{linkstudio/phy/optical.py}{866}{869}
\snipdef{fso_params}{linkstudio/phy/optical.py}{814}{817}
\snipdef{rin}{linkstudio/phy/optical.py}{59}{61}
\snipdef{mmf_delay}{linkstudio/phy/multimode.py}{63}{66}
\snipdef{mmf_rms}{linkstudio/phy/multimode.py}{74}{80}
\snipdef{mmf_chromatic}{linkstudio/phy/multimode.py}{94}{96}
\snipdef{mmf_groups}{linkstudio/phy/multimode.py}{194}{208}
\snipdef{mmf_launch}{linkstudio/phy/multimode.py}{108}{125}
\snipdef{mmf_response}{linkstudio/phy/multimode.py}{237}{243}
\snipdef{mmf_default_alpha}{linkstudio/phy/multimode.py}{184}{190}
\snipdef{laser_rhs}{linkstudio/phy/lasers.py}{142}{149}
\snipdef{laser_steady}{linkstudio/phy/lasers.py}{105}{109}
\snipdef{laser_small_signal}{linkstudio/phy/lasers.py}{236}{243}
\snipdef{laser_overlap}{linkstudio/phy/lasers.py}{90}{99}
\snipdef{vcsel_thermal}{linkstudio/phy/lasers.py}{280}{292}
\snipdef{vcsel_vi}{linkstudio/phy/lasers.py}{267}{270}
\snipdef{laser_block}{linkstudio/phy/lasers.py}{337}{345}
\snipdef{mzm_field}{linkstudio/phy/optical.py}{128}{132}
\snipdef{mzm_block}{linkstudio/phy/optical.py}{229}{233}
\snipdef{raman}{linkstudio/phy/optical.py}{411}{415}
\snipdef{amp_sat}{linkstudio/phy/optical.py}{541}{543}
\snipdef{link_gn}{linkstudio/phy/optical.py}{709}{716}
\snipdef{gn_eta}{linkstudio/phy/theory.py}{233}{240}
\snipdef{osnr_link}{linkstudio/phy/theory.py}{202}{207}
\snipdef{pam_ser}{linkstudio/phy/theory.py}{148}{150}
\snipdef{coh_isnr}{linkstudio/phy/optical.py}{1067}{1072}
\snipdef{delay_kernel}{linkstudio/phy/channel/fading.py}{51}{60}
\snipdef{tr38901}{linkstudio/phy/channel/propagation.py}{35}{56}
\snipdef{dielectric}{linkstudio/phy/wireline.py}{268}{281}
\snipdef{pic}{linkstudio/phy/wireline.py}{248}{257}
\snipdef{fso_channel}{linkstudio/phy/optical.py}{915}{935}

\def\numblocks{90}

\def\numscenarios{42}

\def\dmtrateA{5.2}
\def\dmtrateB{1.3}
\def\dmtrateC{0.3}

\def\fiberLD{1.153}
\def\fiberbeta{21.68}

\def\solitonPzero{0.667}
\def\solitonL{4.61}

\newcommand{\mmfalphaopt}{1.9777}
\newcommand{\mmfalphaomfour}{1.9971}
\newcommand{\mmfalphaomthree}{2.0235}
\newcommand{\mmfdelta}{0.0093}
\newcommand{\mmfgroups}{18}
\newcommand{\mmfsigmamin}{12.2}


% ------------------------------------------------------------- boxes
\newtcolorbox{incode}[1][]{enhanced,breakable,colback=soft,colframe=accent,boxrule=0pt,
  leftrule=2.5pt,arc=0pt,outer arc=0pt,left=8pt,right=8pt,top=5pt,bottom=5pt,
  fonttitle=\sffamily\bfseries\small,coltitle=accent,colbacktitle=soft,
  title={In the code#1},attach title to upper={\quad},before skip=8pt,after skip=8pt}
\newtcolorbox{keypoint}[1][Key idea]{enhanced,breakable,colback=softgreen,colframe=accent3,boxrule=0pt,
  leftrule=2.5pt,arc=0pt,outer arc=0pt,left=8pt,right=8pt,top=5pt,bottom=5pt,
  fonttitle=\sffamily\bfseries\small,coltitle=accent3,colbacktitle=softgreen,
  title={#1},attach title to upper={\quad},before skip=8pt,after skip=8pt}
\newtcolorbox{pitfall}[1][Watch out]{enhanced,breakable,colback=softwarm,colframe=accent2,boxrule=0pt,
  leftrule=2.5pt,arc=0pt,outer arc=0pt,left=8pt,right=8pt,top=5pt,bottom=5pt,
  fonttitle=\sffamily\bfseries\small,coltitle=accent2,colbacktitle=softwarm,
  title={#1},attach title to upper={\quad},before skip=8pt,after skip=8pt}

% ------------------------------------------------------------- inline names
% Inline code names use the url machinery on detokenized text: special characters print
% verbatim, long names can break after / _ . = , and the commands are robust (usable in
% captions). Braces must be balanced.
\makeatletter
\DeclareUrlCommand\cs@tt{\urlstyle{tt}}
\g@addto@macro\UrlBreaks{\do\_\do\.\do\=\do\,}
\DeclareRobustCommand{\code}[1]{\expandafter\cs@tt\expandafter{\detokenize{#1}}}
\DeclareRobustCommand{\file}[1]{\expandafter\cs@tt\expandafter{\detokenize{#1}}}
\DeclareRobustCommand{\blk}[1]{\textcolor{accent}{\code{#1}}}          % block type
\DeclareRobustCommand{\prm}[1]{\textcolor{accent2!85!black}{\code{#1}}}% parameter
\DeclareRobustCommand{\meta}[1]{\textcolor{accent4}{\code{#1}}}        % metadata key
\makeatother
\newcommand{\kind}[1]{\textsf{\small #1}}

% ------------------------------------------------------------- maths
\newcommand{\E}{\mathbb{E}}
\newcommand{\R}{\mathbb{R}}
\newcommand{\C}{\mathbb{C}}
\newcommand{\CN}{\mathcal{CN}}
\newcommand{\N}{\mathcal{N}}
\newcommand{\Q}{\operatorname{Q}}
\newcommand{\erfc}{\operatorname{erfc}}
\newcommand{\sinc}{\operatorname{sinc}}
\newcommand{\diag}{\operatorname{diag}}
\newcommand{\sgn}{\operatorname{sgn}}
\newcommand{\Var}{\operatorname{Var}}
\newcommand{\F}{\mathcal{F}}
\newcommand{\EsNo}{E_s/N_0}
\newcommand{\EbNo}{E_b/N_0}
\newcommand{\dB}{\,\mathrm{dB}}
\newcommand{\T}{^{\mathsf{T}}}
\newcommand{\Her}{^{\mathsf{H}}}
\newcommand{\conj}{^{*}}
\newcommand{\vect}[1]{\bm{#1}}
\newcommand{\mat}[1]{\mathbf{#1}}
\DeclareMathOperator*{\argmin}{arg\,min}
\DeclareMathOperator*{\argmax}{arg\,max}

% ------------------------------------------------------------- TikZ styles
\tikzset{
  >={Latex[length=2mm]},
  blockbox/.style={draw=accent,fill=soft,rounded corners=2pt,minimum height=8mm,minimum width=16mm,
    align=center,font=\sffamily\footnotesize,inner sep=3pt},
  warmbox/.style={blockbox,draw=accent2,fill=softwarm},
  greenbox/.style={blockbox,draw=accent3,fill=softgreen},
  greybox/.style={blockbox,draw=black!40,fill=black!4},
  sumnode/.style={draw=accent,fill=white,circle,inner sep=0pt,minimum size=5mm,font=\footnotesize},
  mulnode/.style={draw=accent,fill=white,circle,inner sep=0pt,minimum size=5mm,font=\footnotesize,
    path picture={\draw[accent] (path picture bounding box.south east)--(path picture bounding box.north west)
    (path picture bounding box.south west)--(path picture bounding box.north east);}},
  sig/.style={->,thick,color=ink},
  lbl/.style={font=\scriptsize\sffamily,color=ink},
  note/.style={font=\scriptsize\itshape,color=commentcolor,align=center},
  layer/.style={draw=#1,fill=#1!7,rounded corners=4pt,inner sep=6pt},
  dot/.style={circle,fill=ink,inner sep=0pt,minimum size=3pt},
}

% ------------------------------------------------------------- pgfplots styles
\pgfplotscreateplotcyclelist{linkstudio}{
  {accent,thick},{accent2,thick},{accent3,thick},{accent4,thick},{accent5,thick},{black!60,thick}}
\pgfplotsset{
  every axis/.append style={font=\footnotesize,
    grid=both, major grid style={black!14}, minor grid style={black!6},
    legend style={font=\scriptsize,draw=black!20,fill=white,fill opacity=0.9,text opacity=1},
    tick label style={font=\scriptsize}, label style={font=\footnotesize},
    cycle list name=linkstudio, unbounded coords=jump, axis line style={black!60}},
  sim/.style={only marks,mark=*,mark size=1.6pt,#1},
  th/.style={no marks,thick,#1},
  simline/.style={thick,mark=*,mark size=1.5pt,#1},
  scatterpts/.style={only marks,mark=*,mark size=0.35pt,#1},
  berplot/.style={xlabel={$E_b/N_0$ [dB]},ylabel={BER},
    yticklabel style={font=\scriptsize}, legend pos=south west},
  constaxis/.style={width=4.4cm,height=4.4cm,axis equal image,xmin=-1.6,xmax=1.6,ymin=-1.6,ymax=1.6,
    xtick={-1,0,1},ytick={-1,0,1},grid=major,title style={font=\scriptsize\sffamily}},
}
\newcommand{\datafile}[1]{generated/data/#1.dat}

% ------------------------------------------------------------- misc
\setcounter{tocdepth}{1}
\setcounter{secnumdepth}{2}
\renewcommand{\arraystretch}{1.15}
\hypersetup{pdftitle={LinkStudio - Theory and User Manual},pdfauthor={LinkStudio project}}
